\documentclass{beamer}
\usepackage[utf8]{inputenc}
\usepackage[vietnamese]{babel}
\usepackage{graphicx}
\usepackage{booktabs}

\usetheme{Madrid}
\usecolortheme{default}

\title[Tiến độ FollowUpAgent]{Báo cáo Tiến độ Dự án FollowUpAgent}
\subtitle{Hoàn thiện Phase 2: AI Pipeline Core}
\author[]{Team UC-04}
\institute{}
\date{23 Tháng 6, 2026}

\begin{document}

\frame{\titlepage}

\begin{frame}
\frametitle{Tóm tắt tiến độ (Executive Summary)}
\begin{itemize}
    \item \textbf{Giai đoạn hoàn thành:} Đã xây dựng hoàn chỉnh \textbf{Phase 2 - AI Pipeline Core}.
    \item \textbf{Các hạng mục kỹ thuật đã triển khai thành công:}
    \begin{itemize}
        \item Xây dựng kiến trúc Database và cấu hình an toàn (Pydantic BaseSettings).
        \item Tích hợp Whisper STT tạo Pipeline tải và bóc băng âm thanh tự động.
        \item Xây dựng dịch vụ AI Local (LLM GGUF) kèm bộ Prompt chấm điểm ứng viên.
        \item Thiết lập dịch vụ Webhook trả kết quả về CRM với cơ chế chống lỗi (Retry/Backoff).
    \end{itemize}
\end{itemize}
\end{frame}

\begin{frame}
\frametitle{Kiến trúc Hệ thống \& Cơ sở dữ liệu}
\begin{itemize}
    \item \textbf{Tự động hoá cấu hình:} Áp dụng Pydantic BaseSettings v2 để quản lý biến môi trường an toàn.
    \item \textbf{Khởi chạy tự động:} Cơ chế tạo bảng Database và nạp mô hình AI ngay khi khởi chạy server.
    \item \textbf{Cơ sở dữ liệu (PostgreSQL/SQLite):} 
    \begin{itemize}
        \item Bảng \texttt{CallRecord}: Lưu trữ thông tin và kết quả phân tích ghi âm phỏng vấn.
        \item Bảng \texttt{FollowUpTask}: Quản lý lộ trình và lập lịch gửi tin nhắn follow-up ứng viên tự động.
    \end{itemize}
\end{itemize}
\end{frame}

\begin{frame}
\frametitle{Pipeline Xử lý Âm thanh (STT)}
\begin{block}{Tự động hoá với PhoWhisper}
Hệ thống xử lý âm thanh tự động qua 3 bước:
\begin{enumerate}
    \item \textbf{Tải file (Download):} Nhận URL từ CRM, tải bằng stream và kiểm tra định dạng.
    \item \textbf{Bóc băng (Transcribe):} Chạy mô hình PhoWhisper-medium (tối ưu tiếng Việt) để chuyển đổi âm thanh thành văn bản.
    \item \textbf{Dọn dẹp (Cleanup):} Tự động xóa file tạm sau xử lý nhằm tiết kiệm dung lượng server.
\end{enumerate}
\end{block}
\end{frame}

\begin{frame}
\frametitle{Pipeline Phân tích AI (Local LLM)}
\begin{itemize}
    \item \textbf{Kiến trúc an toàn:} Sử dụng mô hình GGUF cục bộ (Local LLM), không gửi dữ liệu ra ngoài internet, đảm bảo tuyệt đối tính bảo mật.
    \item \textbf{Đánh giá Ứng viên (HR):} Tích hợp bộ Prompt Template chuyên biệt phục vụ tuyển dụng để:
    \begin{itemize}
        \item Tóm tắt nội dung phỏng vấn, đánh giá năng lực \& thái độ ứng viên.
        \item Tự động tạo gợi ý tin nhắn/email thông báo kết quả.
    \end{itemize}
    \item \textbf{Trích xuất dữ liệu:} Tự động bóc tách và chuẩn hoá phản hồi của AI thành định dạng JSON để đồng bộ.
\end{itemize}
\end{frame}

\begin{frame}
\frametitle{Dịch vụ Giao tiếp CRM (Webhook)}
\begin{itemize}
    \item \textbf{Giao tiếp một đầu mối:} Thiết kế chuyên biệt để nhận file âm thanh từ CRM và trả kết quả về duy nhất hệ thống CRM.
    \item \textbf{Webhook Bất đồng bộ (Asynchronous):}
    \begin{itemize}
        \item Xử lý ngầm tiến trình AI (Nhận audio $\rightarrow$ STT $\rightarrow$ Phân tích LLM).
        \item Tự động đóng gói kết quả tóm tắt và lời khuyên (suggestions) đẩy về endpoint của CRM.
    \end{itemize}
    \item \textbf{Đảm bảo dữ liệu (Reliability):} Sử dụng phương thức HTTP POST với cơ chế \textit{Exponential Backoff} (Tự động thử lại nhiều lần nếu kết nối với CRM bị lỗi).
\end{itemize}
\end{frame}

\begin{frame}
\frametitle{Kế hoạch Tiếp theo (Phase 3: Rule Engine \& Intelligence Layer)}
\begin{itemize}
    \item \textbf{Kết nối Database Rule:} Thiết lập lịch follow-up tự động theo chu kỳ (D+1, D+3, D+7).
    \item \textbf{Map Intent sang Rule:} Ánh xạ ý định (Intent) của ứng viên được AI phân tích vào các quy tắc xử lý tương ứng.
    \item \textbf{Truy vấn ra quyết định:} Khai thác Database để AI tự động quyết định ngày hẹn và luồng kịch bản tiếp theo.
    \item \textbf{Tạo Audit Log:} Xây dựng hệ thống lưu vết (tracking) toàn bộ quá trình ra quyết định của AI để dễ dàng kiểm toán.
\end{itemize}
\end{frame}



\end{document}
